% ===================================================================
%  BAGAN No. 5 — Rumah dan Pembangunannya.
%
%  Traduction, faite en 2026, en indonesien d'aujourd'hui, sur l'ido
%  de texto/io. Regles de la colonne : voir 10-bagan-01.tex. Un seul
%  numero d'alinea, « 01 », comme du cote ido : l'auteur ne numerote
%  aucune replique de ce dialogue.
%
%  SEIZE METIERS, UN SEUL MOT. La ou l'ido construit « -isto » sur
%  une racine — masonisto, tektisto, karpentisto, vitristo,
%  seruristo —, l'indonesien pose « tukang » devant le materiau ou le
%  geste : tukang gali (81) celui qui creuse, tukang batu (108) la
%  pierre, tukang atap (118) le toit, tukang kayu (115) le bois,
%  tukang seng (121) le zinc, tukang plester (99) le platre, tukang
%  cat (102) la peinture, tukang kaca (96) le verre, tukang kunci
%  (97) la serrure, tukang listrik (124) l'electricite, tukang kebun
%  (129) le jardin, tukang kuda (133) le cheval, tukang aduk (87) le
%  gachis, tukang gorden (107) les tentures, tukang timah le plomb,
%  tukang kayu halus le menuisier. Seize fois le meme moule, et pas
%  un emprunt pour nommer l'homme — seulement, parfois, pour nommer
%  sa matiere.
%
%  ET C'EST ICI QU'APPARAIT LA QUATRIEME COUCHE DE CETTE LANGUE. On
%  avait dit trois — le malais de fond, le neerlandais de l'ecole
%  coloniale, l'arabe et le sanskrit du savant. Il en manque une, et
%  ce tableau la met en pleine lumiere : le HOKKIEN, apporte par les
%  marchands et les artisans chinois. « cat » (103), la peinture,
%  vient de 漆 ; « loteng » (34), le grenier, vient de 樓頂, « le haut
%  de l'etage ». Ce sont deux mots du batiment et du commerce urbain,
%  et aucun locuteur ne les sent etrangers. L'en-tete du tableau 1
%  disait trois couches : il en faut dire quatre.
%
%  LE NEERLANDAIS TIENT LES MATERIAUX INDUSTRIELS ET RIEN D'AUTRE :
%  seng le zinc (24), ember le seau (114), pondasi les fondations
%  (2), balkon (28), garasi (51), arsitek (75). Le malais tient tout
%  ce qui se coupe, se porte et se pose : batu, kayu, atap, pasir,
%  kapur, air, tembok, pintu, jendela, tangga, palu, gergaji, kapak,
%  pahat, kikir, kuas, karung, ember mis a part. Le partage suit
%  encore la meme ligne qu'aux tableaux 2, 3 et 4 : ce qu'on manie
%  est malais, ce qui vient d'une usine est neerlandais.
%
%  TROIS FAUTES AU PREMIER JET, ET DEUX SONT DE MISE EN PAGE. Le
%  gipso (70) etait passe derriere les plafonds (26) — l'indonesien
%  place volontiers l'instrument apres le verbe, et il a fallu le
%  remonter en tete de phrase, « Dengan plester ia melapisi... ». Les
%  deux autres sont des fautes de frappe que kolonoj.py a prises
%  seul : un \VUgras{} vide, et un « langit-langit » coupe par la fin
%  de ligne sur son propre trait d'union. Ce dernier est propre a
%  cette colonne : l'indonesien redouble ses mots — langit-langit,
%  unting-unting, anak-anak, bermacam-macam — et le trait d'union y
%  est INTERIEUR au mot, non une cesure. Le controle du trait d'union
%  en fin de ligne, ecrit pour les colonnes europeennes, sert donc ici
%  deux fois plus souvent.
%
%  unting-unting (109), LE FIL A PLOMB, EST LA TROUVAILLE DE CE
%  TABLEAU. Le mot est malais de fond, il nomme exactement l'outil —
%  le poids suspendu qui donne la verticale — et il est redouble
%  comme beaucoup de noms d'objets qui pendent ou oscillent. Une
%  langue qui emprunte « pondasi » au neerlandais garde son propre
%  mot pour l'instrument qui verifie que le mur est droit : c'est
%  l'inverse de ce qu'on attendrait, et c'est la meme logique qu'au
%  tableau 2 — ce qu'on tient dans la main ne s'emprunte pas.
% ===================================================================

%%K t05-apar-1 apar
\VUsaut{6.0mm}
\VUtitre{15.0pt}{60}{BAGAN No. 5}
\VUsaut{1.4mm}
\VUtitre{11.4pt}{0}{Rumah dan Pembangunannya.}
\VUsaut{2.2mm}

%%K t05-tit-1 sub
\VUpk{10.2pt}{40}{Perjumpaan yang baik.}
\VUsaut{3.0mm}

%%K t05-01-1 p
\textsc{Ioannes}. --- Paman, saya ingin sekali tahu bagaimana orang
membangun \VUgras{rumah}.

%%K t05-01-2 p
\textsc{Paman} (80). --- Itu sangat mudah, kawanku ; mari ikut saya,
saya akan menunjukkan kepadamu rumah yang sedang dibangun oleh tuan
Bernardus (79) dan yang akan saya tempati.

%%K t05-01-3 p
\textsc{Ioannes}. --- Siapa yang menggambar \VUgras{denah}
\textsuperscript{(76)} rumah ini?

%%K t05-01-4 p
\textsc{Paman}. --- \VUgras{Arsitek} \textsuperscript{(75)} ; ia
menyerahkan gambar yang sangat jelas dan gambar itu disetujui, lalu
\VUgras{pemborong} \textsuperscript{(77)} memulai pekerjaannya.

%%K t05-01-5 p
\textsc{Ioannes}. --- Lalu siapa pemborong itu?

%%K t05-01-6 p
\textsc{Paman}. --- Ia tuan Blanko, ayah kawanmu Iakobus. Ia memimpin
para pekerja bersama tuan Rufo, sang arsitek.

%%K t05-01-7 p
Nah! Inilah tuan Blanko datang tepat pada waktunya dengan
\VUgras{penggarisnya} \textsuperscript{(78)}, dan tuan Rufo dengan
denahnya.

%%K t05-01-8 p
Selamat pagi, tuan-tuan. Bagaimana kabar Anda?

%%K t05-01-9 p
\textsc{Tuan Rufo}. --- Sangat baik, terima kasih. Selamat pagi,
Ioannes ; alangkah besarnya anak ini sekarang!

%%K t05-01-10 p
\textsc{Paman}. --- Ya, sungguh, ia sudah seperti lelaki kecil. Ia
sangat bersungguh-sungguh ; segala yang dilihatnya harus dijelaskan
kepadanya. Hari ini ia ingin tahu bagaimana orang membangun rumah.

%%K t05-tit-2 sub
\VUpk{10.2pt}{40}{Para pekerja.}
\VUsaut{3.0mm}

%%K t05-01-11 p
\textsc{Tuan Blanko}. --- Nah, ikutlah dengan saya, kawan kecil ; saya
akan segera menjelaskannya kepadamu, sebab saya sangat menyukai
anak-anak yang ingin belajar.

%%K t05-01-12 p
Kau lihat pekerja yang memegang \VUgras{sekop} \textsuperscript{(82)}
di tangan itu : ia \VUgras{tukang gali} \textsuperscript{(81)}. Ia
menggali \VUgras{parit} \textsuperscript{(46)} untuk memasang pipa air.
Ia juga menggali tanah di tempat rumah itu berdiri, supaya tukang batu
dapat mendirikan keempat \VUgras{temboknya}. Bagian tembok yang berada
di dalam tanah disebut \VUgras{pondasi} \textsuperscript{(2)}. Ruang
yang terkurung di antara pondasi itu membentuk \VUgras{ruang bawah
tanah} \textsuperscript{(3)}. Ruang bawah tanah itu mempunyai jendela
kecil yang disebut \VUgras{jendela ruang bawah tanah}
\textsuperscript{(5)}, sebuah \VUgras{pintu} \textsuperscript{(4)} dan
sebuah tangga.

%%K t05-01-13 p
\VUgras{Tukang batu} \textsuperscript{(108)} meneruskan membangun
tembok itu sambil meninggalkan lubang untuk \VUgras{pintu-pintu}
\textsuperscript{(8)} dan \VUgras{jendela-jendela} \textsuperscript{(17)}.

%%K t05-01-14 p
\textsc{Ioannes}. --- Tetapi tukang batu itu tidak saya lihat!

%%K t05-01-15 p
T\textsuperscript{n} \textsc{Blanko}. --- Sekarang ia sudah selesai
bekerja di rumah itu. Ia ada di dekat \VUgras{kandang kuda}
\textsuperscript{(54)}, di atas \VUgras{perancahnya}
\textsuperscript{(110)} ; ia membangun tembok \VUgras{tempat cuci}
\textsuperscript{(56)}.

%%K t05-01-16 p
Ia mempunyai sendok semen, bak adukan dan \VUgras{unting-unting}
\textsuperscript{(109)}. Tetapi nama-nama itu tidak akan kauingat.

%%K t05-01-17 p
\textsc{Ioannes}. --- Tentu saja ingat, sebab saya akan segera meminta
kepada paman saya sebuah sendok semen kecil dan sebuah bak, dan
unting-unting akan saya buat sendiri dengan seutas tali.

%%K t05-tit-3 sub
\VUsaut{3.0mm}
\VUpk{10.2pt}{40}{Bahan bangunan.}
\VUsaut{3.0mm}

%%K t05-01-18 p
\textsc{Paman}. --- Ha! bagus sekali. Tetapi katakanlah kepada saya,
Ioannes, apa yang diperlukan untuk membangun sebuah rumah?

%%K t05-01-19 p
\textsc{Ioannes}. --- Diperlukan \VUgras{batu tembok}
\textsuperscript{(59)} dan \VUgras{adukan} \textsuperscript{(65)}.

%%K t05-01-20 p
\textsc{Paman}. --- Tepat sekali ; lihatlah orang bertopi besar di
atas \VUgras{gerobaknya} \textsuperscript{(136)} itu. Ia \VUgras{kusir
gerobak} \textsuperscript{(135)}. Setiap hari ia membawa
\VUgras{bongkah-bongkah batu} \textsuperscript{(60)} ke sini. Ia
menurunkan muatannya di halaman, dan \VUgras{para pemahat batu}
\textsuperscript{(83)} memotong bongkah itu dan menggergajinya dengan
\VUgras{gergaji batu} \textsuperscript{(85)}. Seorang \VUgras{kuli}
\textsuperscript{(146)} membawanya kepada tukang batu, yang
memasangnya.

%%K t05-01-21 p
Sementara itu, \VUgras{tukang aduk} \textsuperscript{(87)} menyiapkan
adukan. Ia memerlukan \VUgras{pasir} \textsuperscript{(62)},
\VUgras{kapur} \textsuperscript{(63)} dan \VUgras{air}
\textsuperscript{(64)}. Kapur ada di dekatnya dalam \VUgras{karung}
\textsuperscript{(71)} ; seorang \VUgras{pembantu tukang batu}
\textsuperscript{(111)} membawakan pasir kepadanya di atas
\VUgras{gerobak dorong} \textsuperscript{(112)}, dan \VUgras{anak
magang} \textsuperscript{(113)} berdiri di pompa (58). Ia mengisi
\VUgras{ember} \textsuperscript{(114)}. Adukan itu segera siap.
\VUgras{Pengangkut adukan} \textsuperscript{(107)} mengambilnya dan
membawanya naik, lewat tangga, ke atas perancah, tempat seorang tukang
batu menyelesaikan sisi rumah itu.

%%K t05-01-22 p
Dan sekarang, apa sebutan pekerja yang mengerjakan \VUgras{atap}
\textsuperscript{(31)}?

%%K t05-01-23 p
\textsc{Ioannes}. --- Ia \VUgras{tukang atap} \textsuperscript{(118)}.

%%K t05-tit-4 sub
\VUsaut{3.0mm}
\VUpk{10.2pt}{40}{Atap rumah.}
\VUsaut{3.0mm}

%%K t05-01-24 p
T\textsuperscript{n} \textsc{Blanko}. --- Ya, tetapi yang datang lebih
dahulu adalah \VUgras{tukang kayu} \textsuperscript{(115)}.

%%K t05-01-25 p
Tukang kayu mengerjakan \VUgras{rangka kayu} \textsuperscript{(35)}. Ia
menyambung \VUgras{balok-balok} \textsuperscript{(37)} dengan
\VUgras{pasak} \textsuperscript{(117)}. Para pekerja di atas sana,
dengan celana beledu mereka yang lebar, adalah tukang kayu. Yang satu
memegang balok kecil, yang lain memotong pasak dengan \VUgras{kapaknya}
\textsuperscript{(116)}. Yang berada di dekat \VUgras{cerobong asap}
\textsuperscript{(41)} adalah tukang atap. Ia memaku reng-reng di atas
(*) balok, dan \VUgras{lempeng batu tulis} \textsuperscript{(66)} di
atas reng. Kadang-kadang ia memasang genting.

%%K t05-noto-1 noto
\VUnotes{143.2mm}{%
(*) \VUgras{Balk-o} = J. \textit{Balken}, I. \textit{joist}, P.
\textit{solive}, It. \textit{travicello}, R. \textit{balka}, Sp. Port.
\textit{viga}.
Akar teknis yang sampai sekarang belum diterima, tetapi patut
diusulkan untuk menerjemahkan makna kata P. \textit{solive}, yang
bukan trabo (P. \textit{poutre}) dan bukan pula trabeto. Ia adalah
batang kayu bersegi empat yang menopang lantai papan dan
langit-langit.%
}

%%K t05-01-26 p
Atap kandang kuda itu \VUgras{beratap genting} \textsuperscript{(55)},
atap rumah \VUgras{beratap batu tulis} \textsuperscript{(32)}, dan atap
beranda dari \VUgras{seng} \textsuperscript{(24)}.

%%K t05-01-27 p
\textsc{Ioannes}. --- Apakah tukang atap juga mengerjakan atap seng?

%%K t05-01-28 p
T\textsuperscript{n} \textsc{Blanko}. --- Tidak ; yang mengerjakannya
adalah \VUgras{tukang seng} \textsuperscript{(121)} atau \VUgras{tukang
timah,} yaitu yang memasang \VUgras{jendela atap}
\textsuperscript{(38)} di atas atap rumah. Ia juga memasang
\VUgras{talang} \textsuperscript{(43)} dan \VUgras{talang atap}
\textsuperscript{(42)}. Ia menyolder seng dan pelat. Ialah yang
memasang \VUgras{penangkal petir} \textsuperscript{(40)} dan
\VUgras{penunjuk arah angin} \textsuperscript{(39)} di atas
\VUgras{tembok segitiga} \textsuperscript{(36)}.

%%K t05-01-29 p
\textsc{Ioannes}. --- Jadi rumah itu sudah selesai?

%%K t05-tit-5 sub
\VUsaut{3.0mm}
\VUpk{10.2pt}{40}{Perkakas.}
\VUsaut{3.0mm}

%%K t05-01-30 p
T\textsuperscript{n} \textsc{Blanko}. --- Belum seluruhnya.
\VUgras{Tukang plester} \textsuperscript{(99)} membangun dengan
\VUgras{batu bata} \textsuperscript{(61)} dinding-dinding yang membagi
rumah itu menjadi bermacam-macam kamar. Dengan \VUgras{plester}
\textsuperscript{(70)} ia melapisi \VUgras{langit-langit}
\textsuperscript{(26)} dan tembok. Sekarang ia mengerjakan langit-langit
\VUgras{beranda} \textsuperscript{(33)}. Seperti tukang batu, ia
mempunyai \VUgras{sendok semen} \textsuperscript{(100)} dan \VUgras{bak
adukan} \textsuperscript{(101)}.

%%K t05-01-31 p
Kemudian tukang kayu halus mengerjakan pintu, jendela, \VUgras{lantai
papan} \textsuperscript{(16)} dan \VUgras{daun jendela}
\textsuperscript{(19)}.

%%K t05-01-32 p
Sekarang ia bekerja di halaman di atas \VUgras{bangku kerjanya}
\textsuperscript{(90)}. Ia mempunyai \VUgras{gergaji}
\textsuperscript{(94)} untuk menggergaji kayu (69), \VUgras{ketam}
\textsuperscript{(91)}, \VUgras{penjepit} \textsuperscript{(92)},
\VUgras{pahat} \textsuperscript{(94 bis)}, \VUgras{jangka}
\textsuperscript{(95)} dan \VUgras{palu} kayu \textsuperscript{(95 bis)}.

%%K t05-01-33 p
\textsc{Ioannes}. --- Ha! ternyata lebih menyenangkan mengerjakan kayu
daripada memasang batu. Paman, belikanlah untuk saya sebuah bangku
kerja, gergaji dan ketam, sebab saya akan segera membuat kandang untuk
Medor.

%%K t05-01-34 p
\textsc{Paman}. --- Baiklah, nak.

%%K t05-01-35 p
\textsc{Ioannes}. --- Alangkah senangnya saya! Dan siapa yang berdiri
di dekat pintu masuk itu?

%%K t05-tit-6 sub
\VUsaut{3.0mm}
\VUpk{10.2pt}{40}{Pengecatan.}
\VUsaut{3.0mm}

%%K t05-01-36 p
T\textsuperscript{n} \textsc{Blanko}. --- Ia \VUgras{tukang cat}
\textsuperscript{(102)}. Ia berdiri di atas tangganya dengan sekaleng
\VUgras{cat} \textsuperscript{(103)} dan \VUgras{kuasnya}
\textsuperscript{(104)}. Ia mengecat kayu pintu masuk supaya lebih
awet.

%%K t05-01-37 p
Ia juga menempelkan kertas dinding di dalam kamar-kamar.

%%K t05-01-38 p
\VUgras{Tukang gorden} \textsuperscript{(107)} yang berada di atas
\VUgras{tangga} \textsuperscript{(106)} di \VUgras{balkon}
\textsuperscript{(28)} sedang memasang \VUgras{tirai}
\textsuperscript{(29)}.

%%K t05-01-39 p
Pekerja yang berdiri di dekat jendela besar itu adalah \VUgras{tukang
kaca} \textsuperscript{(96)}. Ia memasang \VUgras{kaca jendela}
\textsuperscript{(18)} dengan \VUgras{dempul} \textsuperscript{(73)}.
Pekerja yang lain, yang berlutut di dekat pintu ruang bawah tanah,
adalah \VUgras{tukang kunci} \textsuperscript{(97)}. Ia memasang
\VUgras{kunci} \textsuperscript{(6)}. \VUgras{Kikirnya}
\textsuperscript{(98)} ada di dekatnya.

%%K t05-01-40 p
\textsc{Ioannes}. --- Apakah pekerja yang memukul besi dengan
\VUgras{palunya} \textsuperscript{(84)} itu juga tukang kunci?

%%K t05-01-41 p
T\textsuperscript{n} \textsc{Blanko}. --- Lebih tepatnya, ia pandai
besi (125). Ia menempa \VUgras{barang-barang besi}
\textsuperscript{(67)} untuk \VUgras{tukang listrik}
\textsuperscript{(124)} yang berada di atas atap \VUgras{garasi}
\textsuperscript{(51)}. Ia mempunyai \VUgras{perapian tempa}
\textsuperscript{(126)}, \VUgras{landasan} \textsuperscript{(127)} dan
\VUgras{penjepit} \textsuperscript{(128)}.

%%K t05-01-42 p
Tukang listrik memasang \VUgras{kawat} \VUgras{tembaga}
\textsuperscript{(44)} telepon.

%%K t05-tit-7 sub
\VUpk{10.2pt}{40}{Pekerjaan kebun.}
\VUsaut{3.0mm}

%%K t05-01-43 p
\textsc{Paman}. --- Di ujung \VUgras{halaman} \textsuperscript{(45)},
di dekat \VUgras{pagar besi} \textsuperscript{(47)} dan
\VUgras{gerbang} \textsuperscript{(48)}, kau lihat \VUgras{tukang
kebun} \textsuperscript{(129)}. Ia menyiapkan \VUgras{kebun kecil}
\textsuperscript{(49)}. Ia telah mencangkul tanah dengan
\VUgras{sekopnya} \textsuperscript{(131)}, ia menabur benih dan
menggaru dengan \VUgras{garunya} \textsuperscript{(130)}. Kemudian,
dengan \VUgras{penyiramnya} \textsuperscript{(132)}, ia akan mengairi
\VUgras{bedeng-bedeng} \textsuperscript{(150)} dan tanaman itu akan
tumbuh.

%%K t05-01-44 p
\VUgras{Tukang kuda} \textsuperscript{(133)}, yang baru saja merawat
kuda dan memberinya makan, membersihkan \VUgras{kereta}
\textsuperscript{(52)} dengan spons, \VUgras{pompa tangan}
\textsuperscript{(134)} dan seember air. Kemudian ia akan
memasukkannya kembali ke garasi dan menutup pintunya dengan
\VUgras{palang} \textsuperscript{(53)} yang tebal.

%%K t05-01-45 p
Nah, saya kira kau puas ; kau sudah mendapat penjelasan yang cukup.

%%K t05-01-46 p
\textsc{Ioannes}. --- Ya, paman, tetapi saya ingin sekali melihat
bagian dalam rumah itu.

%%K t05-01-47 p
\textsc{Paman}. --- Itu akan terjadi besok. Tuan Rufo akan
menjelaskan denahnya kepadamu dan menunjukkan nama setiap kamar.

%%K t05-tit-8 sub
\VUsaut{3.0mm}
\VUpk{10.2pt}{40}{Tata ruang dalam rumah.}
\VUsaut{2.4mm}

%%K t05-apar-2 apar
{\centering\textit{(Lihat denahnya.)}\par}
\VUsaut{2.4mm}

%%K t05-01-48 p
\textsc{Arsitek}. --- Mari, Ioannes, saya akan segera mengajakmu
menjelajahi bagian-bagian rumah ini.

%%K t05-01-49 p
Mari kita masuk ke \VUgras{lantai dasar} \textsuperscript{(7)}. Inilah
tombol \VUgras{bel} \textsuperscript{(11)}. Kita menaiki ketiga
\VUgras{anak tangga} \textsuperscript{(14)} \VUgras{tangga depan}
\textsuperscript{(9)}, kita melewati \VUgras{ambang}
\textsuperscript{(10)} \VUgras{pintu masuk} \textsuperscript{(8)}, dan
inilah kita berada di kaki \VUgras{tangga} \textsuperscript{(13)}.

%%K t05-01-50 p
\textsc{Ioannes}. --- Ini lorongnya.

%%K t05-01-51 p
\textsc{Arsitek}. --- Bukan, ini \VUgras{ruang depan}
\textsuperscript{(12)}. \VUgras{Lorong} \textsuperscript{(a)} ada di
ujung sana. Di sebelah kanan, inilah \VUgras{ruang tamu}
\textsuperscript{(b)} dan \VUgras{ruang makan} \textsuperscript{(c)} ;
berhadapan dengan ruang makan ada \VUgras{dapur} \textsuperscript{(d)}
dan \VUgras{ruang pelayan} \textsuperscript{(e)}, lalu di depan
\VUgras{ruang kerja} \textsuperscript{(f)}. \VUgras{Beranda}
\textsuperscript{(23)} dengan \VUgras{pintu kacanya}
\textsuperscript{(25)} ada di dekat ruang tamu.

%%K t05-01-52 p
Sekarang kita akan menaiki tangga. Berpeganglah pada
\VUgras{pegangan tangga} \textsuperscript{(15)} dan hitunglah anak
tangganya. Ada empat belas. Inilah \VUgras{selasar}
\textsuperscript{(m)} ; kita berada di \VUgras{lantai}
\textsuperscript{(27)} pertama.

%%K t05-tit-9 sub
\VUsaut{3.0mm}
\VUpk{10.2pt}{40}{Lantai pertama.}
\VUsaut{3.0mm}

%%K t05-01-53 p
Berhadapan dengan tangga ada \VUgras{jamban} \textsuperscript{(20)} dan
\VUgras{kamar rias} \textsuperscript{(j)}. Di sebelah kanan
sebuah \VUgras{kamar tidur} \textsuperscript{(g)} yang indah dengan dua
jendela menghadap \VUgras{muka rumah} \textsuperscript{(1)}. Di sebelah
kiri ada \VUgras{kamar mandi} \textsuperscript{(1)} dan \VUgras{kamar
tamu} \textsuperscript{(i)} dengan \VUgras{ceruk tidur}
\textsuperscript{(h)} yang besar. Di dekat kamar tidur ada
\VUgras{kamar anak} \textsuperscript{(k)}. Kamar itu menghadap
\VUgras{teras} \textsuperscript{(21)} yang luas. Di bawah teras itu ada
jamban yang lain (20), dan pada dinding, inilah \VUgras{lonceng}
\textsuperscript{(22)} yang berbunyi untuk makan malam.

%%K t05-01-54 p
\textsc{Ioannes}. --- Mari kita naik ke \VUgras{lantai dua}.

%%K t05-01-55 p
\textsc{Arsitek}. --- Tidak ada lantai dua. Di bawah atap hanya ada
\VUgras{kamar-kamar loteng} \textsuperscript{(33)} dan loteng (34).
Penjelajahan kita selesai ; kita boleh turun kembali.

%%K t05-01-56 p
\textsc{Ioannes}. --- Terima kasih, tuan ; semua itu sangat menarik.
Saya akan segera menceritakan kepada ayah saya semua yang saya lihat.

\VUsaut{4.0mm}
\VUfilet{22mm}
